\section{Fundamento teórico}
\label{sec:fundamento-teorico}

\subsection{Modelo de bases de datos distribuidas}

Siguiendo a Özsu y Valduriez~\cite{ozsu2011principles}, un sistema de base de datos distribuido
(SBDD) se define como un conjunto de múltiples bases de datos lógicamente relacionadas,
gestionadas por un único sistema y percibidas por el usuario final como una única base. Un SBDD
persigue cuatro objetivos de diseño: \emph{transparencia} (de fragmentación, de replicación y de
ubicación —el usuario no necesita saber dónde ni cómo se almacena físicamente cada fila—),
\emph{autonomía local} de cada nodo, \emph{disponibilidad} ante fallos parciales, y
\emph{escalabilidad horizontal} mediante la adición de nodos en vez de hardware más potente en un
único nodo.

La \emph{fragmentación horizontal} de una relación $R$ produce fragmentos $R_1, R_2, \dots, R_n$
mediante selección sobre un atributo particionador: $R_i = \sigma_{p_i}(R)$, donde $\{p_i\}$ es un
conjunto de predicados mutuamente excluyentes y colectivamente exhaustivos. Esta fragmentación es
correcta si y solo si cumple tres condiciones:

\begin{itemize}[leftmargin=1.5em]
    \item \textbf{Completitud}: todo elemento de $R$ pertenece a al menos un fragmento $R_i$.
    \item \textbf{Reconstrucción}: $R = \bigcup_{i=1}^{n} R_i$, es decir, la unión de los
        fragmentos reproduce exactamente la relación original.
    \item \textbf{Disyunción}: $R_i \cap R_j = \emptyset$ para todo $i \neq j$, es decir, ningún
        elemento pertenece a más de un fragmento.
\end{itemize}

Para una fragmentación por rango sobre un atributo $a$ con dominio totalmente ordenado, estas tres
condiciones se satisfacen definiendo $n$ puntos de corte $c_0 = -\infty < c_1 < \dots < c_{n-1} <
c_n = +\infty$ y cada fragmento como $R_i = \sigma_{c_{i-1} \leq a < c_i}(R)$: la completitud se
sigue de que el dominio de $a$ está cubierto en su totalidad por los intervalos, la disyunción de
que los intervalos son semiabiertos y consecutivos sin solape, y la reconstrucción de que la unión
de predicados semiabiertos consecutivos es la tautología sobre el dominio de $a$. La
Sección~\ref{sec:diseno-esquema} instancia este esquema con $a = \texttt{created\_at}$, $n = 4$ y
verifica explícitamente las tres condiciones para el caso concreto del equipo.

La \emph{replicación} complementa a la fragmentación para lograr disponibilidad: con un factor de
replicación $f \geq 3$, un sistema tolera la pérdida simultánea de hasta
$\lfloor (f-1)/2 \rfloor$ réplicas por rango de datos, siempre que se mantenga el quórum de
escritura $\lceil f/2 \rceil + 1$. Con $f = 3$ (el valor usado en este despliegue), el quórum de
escritura es 2 y la tolerancia es de 1 réplica: el clúster admite la caída de exactamente un nodo
sin perder disponibilidad de escritura, pero no la de dos simultáneamente. Esta aritmética de
quórum es la que se contrasta empíricamente en la Sección~\ref{sec:cluster-verificacion}.

\subsection{Teorema CAP y modelo PACELC}

El teorema CAP~\cite{brewer2012cap} establece que, ante una partición de red ($P$), un sistema
distribuido no puede garantizar simultáneamente consistencia ($C$) y disponibilidad ($A$).
Abadi~\cite{abadi2012consistency} extiende esta caracterización con PACELC, señalando que incluso
en ausencia de partición ($E$, \emph{else}) todo sistema distribuido enfrenta un segundo
\emph{trade-off} entre latencia ($L$) y consistencia ($C$). CockroachDB se posiciona
explícitamente como un sistema PC/EC: prioriza consistencia sobre disponibilidad ante partición, y
consistencia sobre latencia en operación normal, lo que se manifiesta de forma concreta en el
comportamiento observado durante el incidente operativo documentado en la
Sección~\ref{sec:cluster-verificacion} (rechazo de escrituras concurrentes conflictivas en vez de
aceptarlas y arriesgar una lectura inconsistente posterior).

Como resume Kleppmann~\cite{kleppmann2017ddia} en su tratamiento de sistemas de datos distribuidos
modernos, la elección entre consistencia y disponibilidad no es una propiedad binaria del sistema
completo, sino una decisión de diseño que puede —y en sistemas como CockroachDB, debe— tomarse de
forma explícita y documentada para cada tipo de operación.

\subsection{Consenso distribuido: Raft}

El algoritmo Raft~\cite{ongaro2014raft} descompone el problema del consenso distribuido en tres
subproblemas: \emph{elección de líder}, \emph{replicación de logs} y \emph{seguridad}
(\emph{safety}). En CockroachDB, cada rango de datos constituye un grupo Raft independiente con su
propio líder: una escritura se considera comprometida (\emph{committed}) cuando el líder del rango
recibe confirmación de una mayoría estricta de sus réplicas, garantizando durabilidad frente a la
caída de una minoría de nodos. Con el factor de replicación $f=3$ de la
Sección~\ref{sec:fundamento-teorico} (quórum de escritura 2 de 3), esto explica directamente por
qué el clúster tolera la caída de un nodo sin pérdida de disponibilidad de escritura, pero no la
de dos nodos simultáneamente —ya no existe mayoría posible entre el único nodo restante—,
comportamiento que se contrasta con la medición empírica de la
Sección~\ref{sec:cluster-verificacion}.

\subsection{Modelo de programación paralela}

El modelo RDD (\emph{Resilient Distributed Dataset}) de Apache Spark~\cite{zaharia2012rdd} abstrae
una colección distribuida \textbf{inmutable}, con \textbf{linaje} (\emph{lineage}) de
transformaciones registrado en vez de los datos materializados en cada etapa intermedia,
\textbf{evaluación perezosa} (\emph{lazy evaluation}, las transformaciones no se ejecutan hasta que
una acción las requiere) y \textbf{tolerancia a fallos mediante recómputo}: si una partición se
pierde, Spark la reconstruye reaplicando su linaje de transformaciones sobre los datos de origen,
en vez de depender de una copia de respaldo replicada. Esta abstracción es la que sustenta el
pipeline de la Sección~\ref{sec:pipeline-paralelo}, encadenando cinco transformaciones sin
materializar resultados intermedios en disco entre cada una.

La ganancia de paralelismo del pipeline se interpreta bajo dos modelos complementarios. La ley de
Amdahl~\cite{amdahl1967validity} modela el \emph{speedup} $S(N)$ alcanzable con $N$ unidades de
procesamiento para un problema de \textbf{tamaño fijo}, en función de la fracción $p$ del trabajo
que es paralelizable:

\begin{equation}
    S(N) = \frac{1}{(1-p) + \dfrac{p}{N}}, \qquad \lim_{N \to \infty} S(N) = \frac{1}{1-p}
    \label{eq:amdahl}
\end{equation}

La Ecuación~\ref{eq:amdahl} predice un límite asintótico de \emph{speedup} determinado
enteramente por la fracción secuencial $(1-p)$ del pipeline —en este caso, principalmente la
lectura inicial de Parquet y la escritura final de resultados, que no se benefician de más
particiones más allá de cierto punto—. Gustafson~\cite{gustafson1988reevaluating} señala que esta
formulación asume tamaño de problema constante, y propone en cambio fijar el tiempo total
disponible y dejar crecer el tamaño del problema con $N$; bajo ese supuesto el \emph{speedup}
escala de forma aproximadamente lineal en $N$ en vez de saturar. Ambos modelos se aplican en la
Sección~\ref{sec:protocolo-resultados}: Amdahl para ajustar la fracción paralelizable observada
del pipeline concreto (tamaño de dataset fijo en 600{,}000 filas para todos los niveles de $N$),
y Gustafson como marco de interpretación de por qué, en un caso de uso real de un ISP donde el
volumen de datos crece con el tiempo (no es fijo), el beneficio práctico del paralelismo tiende a
ser mayor que el que predice Amdahl en aislamiento.

\subsection{Diseño y análisis estadístico del experimento}

El contraste de hipótesis del protocolo experimental (Sección~\ref{sec:protocolo-resultados})
sigue el procedimiento estándar para muestras pareadas descrito por Jain~\cite{jain1991art} y
Wohlin et al.~\cite{wohlin2012experimentation}: dado que cada repetición del pipeline paralelo y
del baseline secuencial se ejecuta sobre el mismo dataset y en la misma posición de la secuencia
de repeticiones, las mediciones no son independientes entre grupos, por lo que corresponde una
prueba pareada en vez de una prueba de dos muestras independientes. Antes de aplicar una prueba
paramétrica (\emph{t} de Student pareada) es necesario verificar el supuesto de normalidad de las
diferencias emparejadas; se usa la prueba de Shapiro–Wilk para ello, y en caso de rechazarse la
normalidad ($p \leq 0.05$) se recurre a la alternativa no paramétrica estándar para datos
pareados, la prueba de rangos con signo de Wilcoxon. El intervalo de confianza al 95\% de cada
nivel se calcula como $\bar{t} \pm 1.96 \cdot s/\sqrt{r}$, válido asintóticamente por el teorema
central del límite para $r \geq 8$ muestras por nivel tras el descarte de calentamiento y
enfriamiento descrito en la Sección~\ref{sec:protocolo-resultados}.

\bigskip
\noindent
Los fundamentos anteriores sostienen la capa de datos distribuida heredada de las Entregas 2 y 3.
La Entrega 4 añade sobre ella dos canales de acceso, un pipeline de integración continua y un
protocolo de calidad de producto, cuyo fundamento se desarrolla a continuación.
\label{sec:fundamento-e4}

\subsection{Arquitectura en capas y principios SOLID}

La arquitectura en capas separa un sistema en presentación, aplicación, dominio e
infraestructura~\cite{martin2017clean,evans2003ddd}, con una regla de dependencia única: el
código de una capa solo puede depender de capas más internas, nunca al revés —el dominio no
conoce ni la base de datos ni el protocolo de transporte que lo expone. Los principios SOLID son
las restricciones que impiden que esa separación se erosione con el tiempo: responsabilidad
única (una clase, un motivo de cambio), abierto/cerrado (extensible sin modificar código
existente), sustitución de Liskov (una subclase debe poder reemplazar a su clase base sin alterar
la corrección del programa), segregación de interfaces (interfaces pequeñas y específicas en vez
de una interfaz general que fuerza a implementar métodos irrelevantes), e inversión de
dependencias (los módulos de alto nivel no dependen de los de bajo nivel; ambos dependen de
abstracciones). El uso disciplinado de estos principios es prerrequisito para que los patrones de
diseño del catálogo de Gamma et al.~\cite{gamma1994patterns} aporten valor real, en vez de ser
indirección sin propósito.

\subsection{Ingeniería de aplicaciones móviles}

Las aplicaciones móviles enfrentan retos particulares frente al desarrollo de escritorio o web:
heterogeneidad de dispositivos, restricciones de batería y de red intermitente, un ciclo de vida
gestionado por el sistema operativo (no por el proceso de la propia aplicación), y ventanas de
publicación controladas por tiendas de aplicaciones~\cite{wasserman2010mobile}. La decisión entre
desarrollo nativo y multiplataforma no es una preferencia estética: tiene consecuencias medibles
en rendimiento, mantenibilidad y \emph{time-to-market}, documentadas empíricamente en la
literatura de ingeniería de software móvil~\cite{elkassas2017taxonomy,biornhansen2019crossplatform}.
La Sección~\ref{sec:movil} aplica dos de esos criterios cuantitativos —tamaño de artefacto y
dependencias de interoperabilidad— directamente sobre el proyecto construido, en vez de citarlos
solo como referencia teórica.

\subsection{Integración y entrega continuas}

La entrega continua se apoya en un \emph{pipeline} automatizado que ejecuta pruebas, empaqueta y
despliega en ambientes segregados~\cite{humble2010continuous}. La cultura DevOps subyacente
exige que ningún cambio llegue a un artefacto publicado sin pasar por los mismos
controles~\cite{kim2021devops}: en este proyecto, el pipeline de GitHub Actions
(Sección~\ref{sec:pruebas-cicd}) actúa como \emph{quality gate} explícito mediante dependencias
declaradas entre \emph{jobs} —si la integración falla, el despliegue de imágenes no ocurre.

\subsection{Observabilidad distribuida}

La observabilidad de un sistema distribuido se sustenta en tres señales: métricas (agregados
numéricos con etiquetas), registros estructurados, y trazas (secuencias de \emph{spans}
correlacionados por un identificador común)~\cite{sigelman2010dapper,beyer2016sre}. El estándar
OpenTelemetry~\cite{opentelemetry2024spec} unifica su modelo de datos y su exportación,
permitiendo que distintos proveedores de \emph{backend} de observabilidad (en este proyecto,
Grafana Tempo para trazas y Prometheus para métricas) consuman las mismas señales sin
instrumentación duplicada por proveedor.

\subsection{Modelo de calidad ISO/IEC 25010:2023}

El modelo ISO/IEC 25010:2023~\cite{iso25010} organiza la calidad de producto de software en ocho
características —adecuación funcional, eficiencia de desempeño, compatibilidad, usabilidad,
fiabilidad, seguridad, mantenibilidad y portabilidad—, cada una descompuesta en
subcaracterísticas y evaluable mediante métricas objetivas. Frente a un modelo de calidad
puramente cualitativo, esta estructura obliga a declarar de antemano un umbral numérico por
característica y a reportar la desviación respecto a ese umbral cuando ocurre, en vez de una
valoración subjetiva de si el sistema ``es de buena calidad''. La Sección~\ref{sec:iso25010}
aplica cinco de las ocho características a este proyecto, siguiendo esa misma disciplina de
declarar el umbral antes de medir.
